% GENERATED FILE. Edit canonical lessons, then run npm run generate:books.
% canonical-source-hash: fnv1a64:205b003a0020be15
% canonical-lessons: SA-C277-agatya, SA-C277-prathamam, SA-C277-sighram, SA-C277-mira-nagaram-gatavati, SA-R277-a-day-in-order

\chapter{One Thing, Then the Next}
\label{ch:sa-one-thing-then-the-next}
\hlchaptermodality{277}

\begin{chapteropening}
\textbf{By the end of this chapter:} \emph{I can tell events in order, with -tvā and -ya (gatvā, āgatya) and first, then, in the end (prathamam, tataḥ, ante), and follow a short story told that way.}

Complete SA-R277-a-day-in-order: tell a short day in order, one event after another.
\end{chapteropening}

\section[\texorpdfstring{\sk{आगत्य} (āgatya)}{आगत्य (āgatya)}]{\sk{आगत्य} (āgatya) — one action, then the next}

\label{lesson:SA-C277-agatya}



\begin{quote}
Before the new one: say \textquotedblleft{}the boy, having gone home, reads\textquotedblright{}. (\emph{bālaḥ gṛhaṃ gatvā paṭhati}.) Then say \textquotedblleft{}he went home\textquotedblright{}. (\emph{saḥ gṛhaṃ gatavān}.) Now put the two together.
\end{quote}

\subsection*{You'll want to know: -\sk{त्वा} with the past}
\begin{quote}
\textbf{\sk{सः} \sk{खादित्वा} \sk{गतवान्।}}
\end{quote}

\begin{quote}
\emph{saḥ khāditvā gatavān.} — \textquotedblleft{}He ate, then he left.\textquotedblright{}
\end{quote}

The \textbf{-\sk{त्वा}} form finishes the first action, and the last verb carries the past for both. \textbf{\sk{खादित्वा}} (\emph{khāditvā}) is \textquotedblleft{}having eaten\textquotedblright{}, and \textbf{\sk{कृत्वा}} (\emph{kṛtvā}) \textquotedblleft{}having done\textquotedblright{}:

\begin{quote}
\textbf{\sk{सा} \sk{गृहं} \sk{गत्वा} \sk{पठितवती।}}
\end{quote}

\begin{quote}
\emph{sā gṛhaṃ gatvā paṭhitavatī.} — \textquotedblleft{}She went home and read.\textquotedblright{}
\end{quote}

\begin{quote}
\textbf{\sk{ते} \sk{कार्यं} \sk{कृत्वा} \sk{गतवन्तः।}}
\end{quote}

\begin{quote}
\emph{te kāryaṃ kṛtvā gatavantaḥ.} — \textquotedblleft{}They did the work and left.\textquotedblright{}
\end{quote}

Both actions belong to one doer, and the \textbf{-\sk{त्वा}} action comes first.

\subsection*{You'll want to know: -\sk{य} after a prefix}
A verb with a prefix, such as \textbf{\sk{आगच्छति}} (\emph{āgacchati}), \textquotedblleft{}comes\textquotedblright{}, or \textbf{\sk{प्रविशति}} (\emph{praviśati}), \textquotedblleft{}goes in\textquotedblright{}, takes \textbf{-\sk{य}} (\emph{-ya}) in place of \textbf{-\sk{त्वा}}: one short ending, after a prefix (Pāṇini 7.1.37).

\begin{quote}
\textbf{\sk{सः} \sk{आगत्य} \sk{पठति।}}
\end{quote}

\begin{quote}
\emph{saḥ āgatya paṭhati.} — \textquotedblleft{}He comes and reads.\textquotedblright{}
\end{quote}

\begin{quote}
\textbf{\sk{मीरा} \sk{गृहं} \sk{प्रविश्य} \sk{उपविशति।}}
\end{quote}

\begin{quote}
\emph{mīrā gṛhaṃ praviśya upaviśati.} — \textquotedblleft{}Mira goes into the house and sits down.\textquotedblright{}
\end{quote}

\begin{quote}
\textbf{\sk{सा} \sk{उपविश्य} \sk{खादितवती।}}
\end{quote}

\begin{quote}
\emph{sā upaviśya khāditavatī.} — \textquotedblleft{}She sat down and ate.\textquotedblright{}
\end{quote}

\noindent
\begin{tabularx}{\linewidth}{@{}>{\raggedright\arraybackslash}X>{\raggedright\arraybackslash}X@{}}
\toprule
\textbf{the verb} & \textbf{the action finished} \\
\midrule
\textbf{\sk{गच्छति}} & \textbf{\sk{गत्वा}} (\emph{gatvā}) \\
\textbf{\sk{आगच्छति}} & \textbf{\sk{आगत्य}} (\emph{āgatya}) \\
\textbf{\sk{प्रविशति}} & \textbf{\sk{प्रविश्य}} (\emph{praviśya}) \\
\textbf{\sk{उपविशति}} & \textbf{\sk{उपविश्य}} (\emph{upaviśya}) \\
\bottomrule
\end{tabularx}

\subsection*{Your turn}
Say these aloud:

\begin{itemize}
\raggedright
  \item \textquotedblleft{}Having come, he reads\textquotedblright{} — \emph{saḥ āgatya paṭhati}
  \item \textquotedblleft{}Having come, he eats\textquotedblright{} — \emph{saḥ āgatya khādati}
  \item \textquotedblleft{}Having come, he sits down\textquotedblright{} — \emph{saḥ āgatya upaviśati}
\end{itemize}

\begin{tcolorbox}[breakable,colback=teal!4,colframe=teal!35!black,title={Before you move on}]
\textquotedblleft{}Having gone in\textquotedblright{}? (\textbf{\sk{प्रविश्य}}.) Why \emph{-ya} and not \emph{-tvā}? (The verb has a prefix.) \textquotedblleft{}She ate and then left\textquotedblright{}? (\textbf{\sk{सा} \sk{खादित्वा} \sk{गतवती}}.)
\end{tcolorbox}
\section[\texorpdfstring{\sk{प्रथमम्} · \sk{ततः} · \sk{अन्ते} (prathamam …)}{प्रथमम् · ततः · अन्ते (prathamam …)}]{\sk{प्रथमम्} · \sk{ततः} · \sk{अन्ते} (prathamam · tataḥ · ante) — first, then, in the end}

\label{lesson:SA-C277-prathamam}



\begin{quote}
Say \textquotedblleft{}she went home and read\textquotedblright{}. (\emph{sā gṛhaṃ gatvā paṭhitavatī}.) Two actions in one sentence. A story needs more sentences, and words to order them.
\end{quote}

\subsection*{You'll want to know: \sk{प्रथमम्}, \sk{ततः}, \sk{अनन्तरम्}}
\textbf{\sk{प्रथमम्}} (\emph{prathamam}), \textquotedblleft{}first\textquotedblright{}, is \textbf{\sk{प्रथमः}} (\emph{prathamaḥ}), \textquotedblleft{}first\textquotedblright{}, said as an adverb. \textbf{\sk{ततः}} (\emph{tataḥ}), \textquotedblleft{}then\textquotedblright{}, and \textbf{\sk{अनन्तरम्}} (\emph{anantaram}), \textquotedblleft{}afterwards\textquotedblright{}, you know.

\begin{quote}
\textbf{\sk{प्रथमं} \sk{सः} \sk{खादितवान्।}}
\end{quote}

\begin{quote}
\emph{prathamaṃ saḥ khāditavān.} — \textquotedblleft{}First he ate.\textquotedblright{}
\end{quote}

\begin{quote}
\textbf{\sk{ततः} \sk{सः} \sk{पठितवान्।}}
\end{quote}

\begin{quote}
\emph{tataḥ saḥ paṭhitavān.} — \textquotedblleft{}Then he read.\textquotedblright{}
\end{quote}

\begin{quote}
\textbf{\sk{अनन्तरं} \sk{सः} \sk{गतवान्।}}
\end{quote}

\begin{quote}
\emph{anantaraṃ saḥ gatavān.} — \textquotedblleft{}Afterwards he left.\textquotedblright{}
\end{quote}

\subsection*{You'll want to know: \sk{अन्ते}}
\textbf{\sk{अन्ते}} (\emph{ante}), \textquotedblleft{}in the end\textquotedblright{}, closes the story.

\begin{quote}
\textbf{\sk{अन्ते} \sk{सः} \sk{गृहम्} \sk{आगतवान्।}}
\end{quote}

\begin{quote}
\emph{ante saḥ gṛham āgatavān.} — \textquotedblleft{}In the end he came home.\textquotedblright{}
\end{quote}

Four lines, in order. \textbf{\sk{क्रीतवती}} (\emph{krītavatī}) is the past of \textbf{\sk{क्रीणाति}} (\emph{krīṇāti}), \textquotedblleft{}buys\textquotedblright{}:

\begin{quote}
\textbf{\sk{प्रथमं} \sk{मीरा} \sk{आपणं} \sk{गतवती।}}
\end{quote}

\begin{quote}
\emph{prathamaṃ mīrā āpaṇaṃ gatavatī.} — \textquotedblleft{}First Mira went to the shop.\textquotedblright{}
\end{quote}

\begin{quote}
\textbf{\sk{ततः} \sk{सा} \sk{फलं} \sk{क्रीतवती।}}
\end{quote}

\begin{quote}
\emph{tataḥ sā phalaṃ krītavatī.} — \textquotedblleft{}Then she bought a fruit.\textquotedblright{}
\end{quote}

\begin{quote}
\textbf{\sk{अनन्तरं} \sk{सा} \sk{गृहम्} \sk{आगतवती।}}
\end{quote}

\begin{quote}
\emph{anantaraṃ sā gṛham āgatavatī.} — \textquotedblleft{}Afterwards she came home.\textquotedblright{}
\end{quote}

\begin{quote}
\textbf{\sk{अन्ते} \sk{सा} \sk{फलं} \sk{खादितवती।}}
\end{quote}

\begin{quote}
\emph{ante sā phalaṃ khāditavatī.} — \textquotedblleft{}In the end she ate the fruit.\textquotedblright{}
\end{quote}

\subsection*{Your turn}
Say these aloud:

\begin{itemize}
\raggedright
  \item \textquotedblleft{}First I ate\textquotedblright{} — \emph{prathamam ahaṃ khāditavān}
  \item \textquotedblleft{}Then I ate\textquotedblright{} — \emph{tataḥ ahaṃ khāditavān}
  \item \textquotedblleft{}In the end I ate\textquotedblright{} — \emph{ante ahaṃ khāditavān}
\end{itemize}

\begin{tcolorbox}[breakable,colback=teal!4,colframe=teal!35!black,title={Before you move on}]
\textquotedblleft{}First\textquotedblright{}? (\textbf{\sk{प्रथमम्}}.) \textquotedblleft{}Afterwards\textquotedblright{}? (\textbf{\sk{अनन्तरम्}}.) \textquotedblleft{}In the end\textquotedblright{}? (\textbf{\sk{अन्ते}}.) Which of them opens a story? (\textbf{\sk{प्रथमम्}}.)
\end{tcolorbox}
\section[\texorpdfstring{\sk{शीघ्रम्} (śīghram)}{शीघ्रम् (śīghram)}]{\sk{शीघ्रम्} (śīghram) — quickly}

\label{lesson:SA-C277-sighram}



\begin{quote}
Before the new one: say \textquotedblleft{}first he ate, then he left\textquotedblright{}. (\emph{prathamaṃ saḥ khāditavān, tataḥ saḥ gatavān}.)
\end{quote}

\subsection*{You'll want to know: \sk{शीघ्रम्}}
\textbf{\sk{शीघ्रम्}} — \emph{śīghram} — \textquotedblleft{}quickly\textquotedblright{}.

\begin{quote}
\textbf{\sk{शीघ्रम्} \sk{आगच्छतु।}}
\end{quote}

\begin{quote}
\emph{śīghram āgacchatu.} — \textquotedblleft{}Please come quickly.\textquotedblright{}
\end{quote}

\begin{quote}
\textbf{\sk{सः} \sk{शीघ्रं} \sk{धावति।}}
\end{quote}

\begin{quote}
\emph{saḥ śīghraṃ dhāvati.} — \textquotedblleft{}He runs quickly.\textquotedblright{}
\end{quote}

\begin{quote}
\textbf{\sk{सा} \sk{शीघ्रं} \sk{गृहं} \sk{गतवती।}}
\end{quote}

\begin{quote}
\emph{sā śīghraṃ gṛhaṃ gatavatī.} — \textquotedblleft{}She went home quickly.\textquotedblright{}
\end{quote}

\begin{cousinweb}
\textbf{\sk{शीघ्रम्}} (\emph{śīghram}) is the neuter of \textbf{\sk{शीघ्र}} (\emph{śīghra}), \textquotedblleft{}quick, swift\textquotedblright{}, said as an adverb, as \textbf{\sk{प्रथमम्}} is. Monier-Williams calls its derivation doubtful, so its origin is uncertain. It lives on as Sindhi \emph{sigho}, \textquotedblleft{}quick\textquotedblright{} (CDIAL 12484).
\end{cousinweb}

\subsection*{Your turn}
Say these aloud:

\begin{itemize}
\raggedright
  \item \emph{śīghram}
  \item \textquotedblleft{}Please come quickly\textquotedblright{} — \emph{śīghram āgacchatu}
  \item \textquotedblleft{}First he came quickly\textquotedblright{} — \emph{prathamaṃ saḥ śīghram āgatavān}
\end{itemize}

\begin{tcolorbox}[breakable,colback=teal!4,colframe=teal!35!black,title={Before you move on}]
What does \textbf{\sk{शीघ्रम्}} mean? (\textquotedblleft{}Quickly\textquotedblright{}.) \textquotedblleft{}She went home quickly\textquotedblright{}? (\textbf{\sk{सा} \sk{शीघ्रं} \sk{गृहं} \sk{गतवती}}.)
\end{tcolorbox}
\section[(reading)]{(reading) — Mira's day in the city}

\label{lesson:SA-C277-mira-nagaram-gatavati}



\begin{quote}
Eighty words, all met before but two new forms: \textbf{\sk{प्रविष्टवती}} (\emph{praviṣṭavatī}), \textquotedblleft{}she went in\textquotedblright{}, and \textbf{\sk{दृष्ट्वा}} (\emph{dṛṣṭvā}), \textquotedblleft{}having seen\textquotedblright{}.
\end{quote}

\begin{tcolorbox}[breakable,skin=enhanced,colback=orange!6,colframe=orange!45!black,boxrule=0.5pt,arc=1mm,left=6pt,right=6pt,top=4pt,bottom=4pt,fonttitle=\bfseries,title={Read this}]
\begin{quote}
\textbf{\sk{ह्यः} \sk{मीरा} \sk{नगरं} \sk{गतवती।}}
\end{quote}

\begin{quote}
\emph{hyaḥ mīrā nagaraṃ gatavatī.}
\end{quote}

\begin{quote}
\textbf{\sk{प्रथमं} \sk{सा} \sk{स्थानकं} \sk{गतवती।}}
\end{quote}

\begin{quote}
\emph{prathamaṃ sā sthānakaṃ gatavatī.}
\end{quote}

\begin{quote}
\textbf{\sk{तत्र} \sk{सा} \sk{एकं} \sk{यानपत्रं} \sk{क्रीतवती।}}
\end{quote}

\begin{quote}
\emph{tatra sā ekaṃ yānapatraṃ krītavatī.}
\end{quote}

\begin{quote}
\textbf{\sk{ततः} \sk{सा} \sk{धूमयानं} \sk{प्रविष्टवती।}}
\end{quote}

\begin{quote}
\emph{tataḥ sā dhūmayānaṃ praviṣṭavatī.}
\end{quote}

\begin{quote}
\textbf{\sk{धूमयाने} \sk{सा} \sk{पुस्तकं} \sk{पठति} \sk{स्म।}}
\end{quote}

\begin{quote}
\emph{dhūmayāne sā pustakaṃ paṭhati sma.}
\end{quote}

\begin{quote}
\textbf{\sk{यदा} \sk{सा} \sk{पठति} \sk{स्म}, \sk{तदा} \sk{वर्षा} \sk{आगतवती।}}
\end{quote}

\begin{quote}
\emph{yadā sā paṭhati sma, tadā varṣā āgatavatī.}
\end{quote}

\begin{quote}
\textbf{\sk{धूमयानं} \sk{शीघ्रं} \sk{गच्छति} \sk{स्म।}}
\end{quote}

\begin{quote}
\emph{dhūmayānaṃ śīghraṃ gacchati sma.}
\end{quote}

\begin{quote}
\textbf{\sk{अनन्तरं} \sk{सा} \sk{नगरम्} \sk{आगतवती।}}
\end{quote}

\begin{quote}
\emph{anantaraṃ sā nagaram āgatavatī.}
\end{quote}

\begin{quote}
\textbf{\sk{नगरे} \sk{सा} \sk{आपणं} \sk{गत्वा} \sk{फलानि} \sk{क्रीतवती।}}
\end{quote}

\begin{quote}
\emph{nagare sā āpaṇaṃ gatvā phalāni krītavatī.}
\end{quote}

\begin{quote}
\textbf{\sk{तत्र} \sk{अरुणः} \sk{अपि} \sk{आगतवान्।}}
\end{quote}

\begin{quote}
\emph{tatra aruṇaḥ api āgatavān.}
\end{quote}

\begin{quote}
\textbf{\sk{मीरा} \sk{अरुणं} \sk{दृष्ट्वा} \sk{प्रसन्ना} \sk{अभवत्।}}
\end{quote}

\begin{quote}
\emph{mīrā aruṇaṃ dṛṣṭvā prasannā abhavat.}
\end{quote}

\begin{quote}
\textbf{\sk{अरुणः} \sk{मीरा} \sk{च} \sk{उद्यानं} \sk{गतवन्तौ।}}
\end{quote}

\begin{quote}
\emph{aruṇaḥ mīrā ca udyānaṃ gatavantau.}
\end{quote}

\begin{quote}
\textbf{\sk{उद्याने} \sk{बालाः} \sk{क्रीडन्ति} \sk{स्म।}}
\end{quote}

\begin{quote}
\emph{udyāne bālāḥ krīḍanti sma.}
\end{quote}

\begin{quote}
\textbf{\sk{अरुणः} \sk{तत्र} \sk{उपविश्य} \sk{फलं} \sk{खादितवान्।}}
\end{quote}

\begin{quote}
\emph{aruṇaḥ tatra upaviśya phalaṃ khāditavān.}
\end{quote}

\begin{quote}
\textbf{\sk{ततः} \sk{वर्षा} \sk{पुनः} \sk{आगतवती।}}
\end{quote}

\begin{quote}
\emph{tataḥ varṣā punaḥ āgatavatī.}
\end{quote}

\begin{quote}
\textbf{\sk{अन्ते} \sk{मीरा} \sk{शीघ्रं} \sk{गृहम्} \sk{आगतवती।}}
\end{quote}

\begin{quote}
\emph{ante mīrā śīghraṃ gṛham āgatavatī.}
\end{quote}

\begin{quote}
\textbf{\sk{श्वः} \sk{सा} \sk{पुनः} \sk{नगरं} \sk{गमिष्यति।}}
\end{quote}

\begin{quote}
\emph{śvaḥ sā punaḥ nagaraṃ gamiṣyati.}
\end{quote}

\emph{Read it:} the passage a second time with the romanization covered, once through, without stopping; do not translate as you go
\end{tcolorbox}

\begin{culture}
How this book writes: Pāṇini (1.4.109) calls words said close together \emph{saṃhitā}, and there the sounds at the joins change. Where this book has \emph{bālaḥ gacchati}, a classical page has \textbf{\sk{बालो} \sk{गच्छति}} (\emph{bālo gacchati}), and \emph{saḥ gacchati} is \textbf{\sk{स} \sk{गच्छति}} (\emph{sa gacchati}; Pāṇini 6.1.132). This book writes each word whole, with the one join you know, \emph{gṛhaṃ gacchati}. Recognise the joined forms; there is nothing here to apply.
\end{culture}

\subsection*{Your turn}
Say these aloud:

\begin{itemize}
\raggedright
  \item what Mira did first, at the station
  \item how the day ended, with \emph{ante}
\end{itemize}

\begin{tcolorbox}[breakable,colback=teal!4,colframe=teal!35!black,title={Before you move on}]
Why \emph{paṭhati sma} in the train? (The reading was going on when the rain came.) Which words keep the steps in order? (\emph{prathamam}, \emph{tataḥ}, \emph{anantaram}, \emph{ante}.)
\end{tcolorbox}
\section[(review)]{Review — a day, in order}

\label{lesson:SA-R277-a-day-in-order}



\begin{quote}
Say \textquotedblleft{}she went home\textquotedblright{} and \textquotedblleft{}when I was reading, Arun came\textquotedblright{}. (\emph{sā gṛhaṃ gatavatī}, \emph{yadā ahaṃ paṭhāmi sma, tadā aruṇaḥ āgatavān}.)
\end{quote}

\subsection*{Your turn}
Say these aloud:

\begin{itemize}
\raggedright
  \item First I went to the station.
  \item Then I bought a ticket and got on the train.
  \item Mira came in and sat down.
  \item Afterwards we ate quickly.
  \item In the end Arun came home.
  \item Tell Mira's day in the city in four lines, with \emph{prathamam}, \emph{tataḥ}, \emph{anantaram} and \emph{ante}.
\end{itemize}

\begin{tcolorbox}[breakable,colback=teal!4,colframe=teal!35!black,title={Before you move on}]
\textquotedblleft{}Having gone in\textquotedblright{}: \emph{-ya} or \emph{-tvā}? (\emph{-ya}, \textbf{\sk{प्रविश्य}}: the verb has a prefix.) \textquotedblleft{}Having eaten\textquotedblright{}? (\textbf{\sk{खादित्वा}}.) Now the other way round: what does \textbf{\sk{शीघ्रम्}} mean? (\textquotedblleft{}Quickly\textquotedblright{}.) And \textbf{\sk{अन्ते}}? (\textquotedblleft{}In the end\textquotedblright{}.)
\end{tcolorbox}
